% !TEX root = ../DoAn.tex
\documentclass[../DoAn.tex]{subfiles}
\begin{document}

\section*{Kết luận}

Đề tài đã hoàn thành việc phân tích, thiết kế, cài đặt, triển khai và kiểm chứng một hệ thống
thương mại điện tử theo kiến trúc microservices với Java 21, Spring Boot 3.5 và Spring Cloud
2025. Hệ thống gồm 13 business microservice, ba Spring Cloud infrastructure service, các
thành phần hạ tầng dữ liệu/messaging/observability, cùng Next.js storefront và Admin Panel
phục vụ demo các luồng nghiệp vụ chính.

Về kiến trúc, hệ thống áp dụng Domain-Driven Design và Bounded Context để phân rã bài toán
thành các service độc lập theo miền nghiệp vụ. Giao tiếp đồng bộ được thực hiện qua
REST/OpenFeign, trong khi giao tiếp bất đồng bộ sử dụng Apache Kafka. Các mẫu thiết kế phân
tán quan trọng đã được hiện thực và kiểm chứng gồm Saga Orchestration cho luồng đặt hàng,
Transactional Outbox cho order/payment event, Idempotent Consumer để chống xử lý trùng lặp,
CQRS-lite cho tìm kiếm sản phẩm bằng Elasticsearch, scheduler/reconciliation cho các tình
huống timeout và tự phục hồi.

Về tính năng, hệ thống đáp ứng các nghiệp vụ cốt lõi của một nền tảng thương mại điện tử:
đăng ký và xác thực người dùng qua Keycloak, quản lý sản phẩm/danh mục/thương hiệu, giỏ hàng,
đặt hàng COD và VNPAY sandbox, voucher, tồn kho, thông báo email, đánh giá sản phẩm, tìm kiếm
toàn văn, nội dung CMS, flash-sale và quản trị hệ thống. Next.js storefront hỗ trợ các luồng
người mua, còn Admin Panel cung cấp các module quản trị dashboard, products, inventory, orders,
users, vouchers, flash-sales, content và reviews.

Điểm nổi bật nhất về kỹ thuật là luồng đặt hàng và flash-sale. Luồng đặt hàng thể hiện đầy đủ
Saga Orchestration, Outbox và Idempotent Consumer để đạt tính nhất quán cuối cùng giữa order,
inventory, voucher, payment và notification. Luồng flash-sale sử dụng Redis Lua atomic script,
buyer set, compensation và ReconciliationScheduler để chống over-sell. Spike test 500 VU cho
kết quả đúng 100 purchase thành công trên 100 slot, 100 đơn confirmed và 0 duplicate buyer,
cho thấy thiết kế đáp ứng mục tiêu thực nghiệm của đồ án.

Về vận hành, hệ thống được đóng gói bằng Docker Compose và triển khai production-like trên AWS
EC2 single-host. Hệ thống sử dụng Caddy reverse proxy, HTTPS và GitHub Actions/GHCR cho quy
trình build/push image. Observability stack gồm Prometheus, Grafana và Zipkin giúp theo dõi metrics,
trace và debug các luồng nghiệp vụ end-to-end. Kết quả kiểm thử cho thấy backend readiness
đạt 9/9 split suites; catalog soak đạt p95 61.68 ms với 0.00\% lỗi; checkout stress đạt p95
160.92 ms với order success 100.00\%; các kịch bản Kill Kafka, Kill Redis và inventory
failure compensation đều đạt trong phạm vi thực nghiệm.

\section*{Hướng phát triển}

Bên cạnh kết quả đạt được, hệ thống vẫn còn một số hướng phát triển quan trọng nếu tiếp tục
nâng cấp sau phạm vi đồ án.

\textbf{Nâng cấp hạ tầng triển khai.} Hướng phát triển tự nhiên là chuyển từ Docker Compose
single-host sang Kubernetes multi-node để có high availability, autoscaling, rolling update
và self-healing tốt hơn. Khi đó có thể bổ sung Service Mesh như Istio hoặc Linkerd để quản lý
mTLS, traffic splitting, retry policy và fault injection ở tầng hạ tầng.

\textbf{Hoàn thiện CI/CD và vận hành production.} Pipeline hiện đã đáp ứng build/push image
và deploy EC2 ở mức production-like. Giai đoạn tiếp theo nên bổ sung automated test gate,
security scanning bằng Trivy/Snyk, quản lý secret chuyên dụng, deployment đa môi trường,
Alertmanager/Grafana Alerting và log aggregation tập trung bằng ELK hoặc Loki.

\textbf{Hoàn thiện frontend và nghiệp vụ quản trị.} Storefront và Admin Panel đã đủ cho demo
DATN, nhưng để tiến gần sản phẩm thương mại cần bổ sung object storage cho upload ảnh/file,
workflow kiểm duyệt review, ban/unban user qua Keycloak Admin API, dashboard analytics chuyên
dụng, tối ưu mobile UX và các nghiệp vụ sau bán như đổi trả, hoàn tiền, vận chuyển.

\textbf{Nâng cao hợp đồng tích hợp và chất lượng kiểm thử.} Module \texttt{common} đang đóng
vai trò shared kernel cho DTO/event schema, phù hợp monorepo đồ án nhưng có thể tạo coupling
khi hệ thống mở rộng. Hướng phát triển là bổ sung schema registry hoặc contract testing cho
Kafka/REST API, thiết lập ngưỡng coverage với JaCoCo, thêm mutation test cho nghiệp vụ quan
trọng và kiểm thử security theo OWASP Top 10.

\textbf{Tối ưu hiệu năng và dữ liệu.} Với quy mô lớn hơn, hệ thống có thể tách PostgreSQL
thành cluster hoặc managed database, dùng read replica cho các read-heavy service, Redis
Cluster cho cache/flash-sale, storage bền vững cho Zipkin, và analytics/read model riêng cho
báo cáo doanh thu. Các hướng như Event Sourcing hoặc gRPC nội bộ chỉ nên cân nhắc khi nhu cầu
audit hoặc latency thực sự vượt quá khả năng của thiết kế hiện tại.

\end{document}
